\scanpage{11}%
morphisme de transition $H_n\handwritten{\longrightarrow}H_m$ induit un isomorphisme de $H_n$ sur le
sous-$A$-$\underline C'$-module $\operatorname{Hom}_A(A_n,H_m)$ des ``éléments annulés par $I^{n+1}$'' . Bien
entendu, sous l'hypothèse sur $\underline C'$ faite dans 2.\deleted{\unsure{8}}\handwritten{9}. et 2.10. on voit aussitôt directement que la catégorie de systèmes inductifs qu'on vient de
décrire, est équivalente à \deleted{celle menti} la catégorie de $A$-$\underline C'$-modules mentionnée dans 2.10.\handwritten{, qui est \unsure{donc} un cas \unsure{particulier} de 2.11.}

\uline{Corollaire} 2.12. Soient $A$ un anneau comm noeth. muni d'un idéal $I$,
$\underline C$ la catégorie des $A$-modules de type fini de support $\handwritten{\subset}V(I)$ i.e. annulés
par une puissance de $I$, $\underline C'$ une catégorie abélienne à limites inductives
filtrantes dénombrables exactes, $(T^i)$ un foncteur cohomologique
contravariant de $\underline C$ dans $\underline C'$, tel que $T^i=0$ pour $i\handwritten{<}i_{\unsure{o|0}}$, $n$ un entier,
$(M_j)$ une famille de $A$-modules de type fini tel que laréunion des
supports des $M_j$ soit égale à $Y=V(I)$. \deleted{Cond} Posons pour tout entier $i$:
\[
H^i=\underset{\handwritten{\longrightarrow}}{\lim}_{\!m}T^i(A_{\deleted{\unsure{n}}\handwritten{m}})\qquad(\text{ où }A_m=A/I^{m+1})\;.
\]
Alors les conditions suivantes sont équivalentes:

(i) On a $T^i=0$ pour $i\handwritten{<}n$.

(ii) On a $H^i=0$ pour $i\handwritten{<}n$

(iii) On a $T^i(M_j)=0$ pour $i\handwritten{<}n$ et tout $j$.

Dém. L'équivalence de (i) et (iii) est mise pour mémoire, et a été
vue dans 2.4. On a (i)$\handwritten{\Rightarrow}$(ii), pour l'implication inverse on
raisonne par récurrence sur $n$, c'est trivial pour $n\handwritten{\leq}i_o$, pour $n$ quelconque
et supposant que c'est fait pour les $n'\handwritten{<}n$, il reste à prouver que $T^{n\handwritten{-1}}=0$,
or $T^{n\handwritten{-1}}$ est exact à gauche, donc en vertu de 2.9. est isomorphe à
$\operatorname{Hom}_A(M,H^{n\handwritten{-1}})$, \inserted{\handwritten{donc nul,}} ce qui achève la démonstration. De plus, on voit:

\uline{Corollaire} 2.13. Si les conditions équivalentes (i) à (iii) de 2.12. sont
vérifiées, on a un isomorphisme fonctoriel
\[
T^{n\deleted{\unsure{-1}}}(M)=\operatorname{Hom}_A(M,H^{n\deleted{\unsure{-1}}})
\]

Il serait préférable d'inclure 2.13. dans l'énoncé de 2.12. et de
démontrer par récurrence (ii)$\handwritten{\Rightarrow}$(i) \uline{plus} l'énoncé final 2.13. , pour évi\deleted{\unsure{t}}ter des red\deleted{\unsure{q}}ites.
% archive sheet number: 10

